\subsection{Decomposition and inertia for integrally closed domains}

LEMMA 4. Let A be an integrally closed domain, K its field of fractions, p the characteristics exponent of K, K' a radicial extension of K and A' a subring of K' containing A and integral over A. For every prime ideal p of A, there exists a unique prime ideal \(\mathfrak{p}'\) of \(\mathbf{A}'\) lying above \(\mathfrak{p}\) and \(\mathfrak{p}'\) is the set of \(x \in \mathbf{A}'\) such that there exists an integer \(m \geqslant 0\) for which \(x^{p^m} \in \mathfrak{p}\) .

The existence of \(\mathfrak{p}'\) follows from no. 1, Theorem 1. If \(x \in \mathfrak{p}'\) , there exists \(m \geqslant 0\) such that \(x^{p^m} \in \mathbf{K}\) , whence \(x^{p^m} \in \mathbf{A}\) since \(\mathbf{A}\) is integrally closed, hence \(x^{p^m} \in \mathfrak{p}' \cap \mathbf{A} = \mathfrak{p}\) . Conversely, if \(x \in \mathbf{A}'\) is such that \(x^{p^m} \in \mathfrak{p} \subset \mathfrak{p}'\) , then \(x \in \mathfrak{p}\) since \(\mathfrak{p}'\) is prime.

Remark (1). It follows from § 1, no. 3, Corollary to Proposition 11 that the integral closure of A in \(\mathbf{K}'\) is the set of \(x \in \mathbf{K}'\) such that there exists \(m \geqslant 0\) for which \(x^{p^m} \in \mathbf{A}\) (Algebra, Chapter V, § 8, no. 1, Proposition 1).

PROPOSITION 6. Let A be an integrally closed domain, K its field of fractions, K' a quasi-Galois extension of K and A' the integral closure of A in K'. Then:

\begin{enumerate}
\def\labelenumi{(\roman{enumi})}
\item
  For every prime ideal p of A, the group G of K-automorphisms of K' operates transitively on the set of prime ideals of A' lying over p.
\item
  For every prime ideal \(\mathfrak{p}'\) of \(\mathbf{A}'\) , the field of fractions \(k'\) of \(\mathbf{A}' / \mathfrak{p}'\) is a quasi-Galois extension of the field of fractions \(k\) of \(\mathbf{A} / (\mathbf{A} \cap \mathfrak{p}')\) and the canonical homomorphism \(\sigma \mapsto \bar{\sigma}\) from \(\mathcal{G}^{\mathbb{Z}}(\mathfrak{p}')\) to the group \(\Gamma\) of \(k\) -automorphisms of \(k'\) defines, by taking the quotient, a bijection of \(\mathcal{G}^{\mathbb{Z}}(\mathfrak{p}') / \mathcal{G}^{\mathbb{T}}(\mathfrak{p}')\) onto \(\Gamma\) .
\end{enumerate}

\begin{enumerate}
\def\labelenumi{(\Alph{enumi})}
\item
  Suppose first that \(\mathbf{K}'\) is a finite Galois extension of \(\mathbf{K}\) . Then \(\mathbf{A} = \mathbf{A}' \cap \mathbf{K}\) since \(\mathbf{A}\) is integrally closed and \(\mathbf{A}\) is therefore the ring of invariants of \(\mathcal{G}\) in \(\mathbf{A}'\) . As \(\mathcal{G}\) is finite, the proposition follows in this case from no. 2, Theorem 2.
\item
  Suppose secondly that \(\mathbf{K}'\) is any Galois extension of \(\mathbf{K}\) . Then \(\mathbf{K}'\) is the union of a right directed family \((\mathrm{K}_{\alpha})_{\alpha \in I}\) of finite Galois extensions of \(\mathbf{K}\) . To show (i), consider two prime ideals \(\mathfrak{p}'\) , \(\mathfrak{q}'\) of \(\mathbf{A}'\) lying above \(\mathfrak{p}\) . For all \(\alpha \in I\) , \(\mathfrak{p}' \cap \mathrm{K}_{\alpha}\) and \(\mathfrak{q}' \cap \mathrm{K}_{\alpha}\) are two prime ideals of \(\mathbf{A}' \cap \mathrm{K}_{\alpha}\) lying above \(\mathfrak{p}\) . Since \(\mathbf{A}' \cap \mathrm{K}_{\alpha}\) is the integral closure of \(A\) in \(\mathrm{K}_{\alpha}\) and the restrictions to \(\mathrm{K}_{\alpha}\) of the elements of \(\mathcal{G}\) form the group of K-automorphisms of \(\mathbf{K}\) , it follows from case (A) that there exists \(\sigma \in \mathcal{G}\) such that \(\sigma. (\mathfrak{p}' \cap \mathrm{K}_{\alpha}) = \mathfrak{q}' \cap \mathrm{K}_{\alpha}\) . Let \(\mathcal{E}_{\alpha}\) be the set of \(\sigma \in \mathcal{G}\) which have the above property. Let \(\sigma \in \mathcal{G} - \mathcal{E}_{\alpha}\) ; then, for all \(\tau \in \mathcal{G}\) leaving invariant the elements of \(\mathrm{K}_{\alpha}, (\sigma \tau). (\mathfrak{p}' \cap \mathrm{K}_{\alpha}) = \sigma. (\mathfrak{p}' \cap \mathrm{K}_{\alpha}) \neq \mathfrak{q}' \cap \mathrm{K}_{\alpha}\) and hence \(\sigma \tau \in \mathcal{G} - \mathcal{E}_{\alpha}\) . It follows that \(\mathcal{E}_{\alpha}\) is closed in the topological Galois group \(\mathcal{G}\) (Algebra, Chapter V, Appendix II, no. 1) and clearly the family \((\mathcal{E}_{\alpha})_{\alpha \in I}\) is left directed. As \(\mathcal{G}\) is compact (loc. cit., no. 2, Proposition 3) and the \(\mathcal{E}_{\alpha}\) are non-empty, the intersection \(\mathcal{E}\) of the family \((\mathcal{E}_{\alpha})\) is non-empty and \(\sigma. \mathfrak{p}' = \mathfrak{q}'\) for all \(\sigma \in \mathcal{E}\) , whence (i).
\end{enumerate}

To show (ii), note that \(k'\) is the union of the right directed family \((k_{\alpha})_{\alpha \in \mathbf{I}}\) , where \(k_{\alpha}\) is the field of fractions of \((\mathrm{A}' \cap \mathrm{K}_{\alpha}) / (\mathfrak{p}' \cap \mathrm{K}_{\alpha})\) . As each \(k_{\alpha}\) is a quasi-Galois extension of \(k\) by (A), so is \(k'\) (Algebra, Chapter V, § 6, no. 3, Proposition 8). On the other hand, let \(u\) be a \(k\) -automorphism of \(k'\) and let \(\pi': \mathrm{A}' \to \mathrm{A}' / \mathfrak{p}'\) be the canonical homomorphism. By virtue of no. 2, Theorem 2 applied to \(\mathrm{A}' \cap \mathrm{K}_{\alpha}\) , there exists for all \(\alpha\) a non-empty set \(\mathcal{F}_{\alpha}\) of elements \(\sigma \in \mathcal{G}\) such that \(\sigma \cdot (\mathfrak{p}' \cap \mathrm{K}_{\alpha}) = \mathfrak{p}' \cap \mathrm{K}_{\alpha}\) and \(u(\pi'(x)) = \pi'(\sigma.x)\) for all \(x \in \mathrm{A}' \cap \mathrm{K}_{\alpha}\) . As above it is seen that \(\mathcal{F}_{\alpha}\) is closed in \(\mathcal{G}\) and, as \((\mathcal{F}_{\alpha})\) is a left directed set, its intersection \(\mathcal{F}\) is non-empty. Clearly for \(\sigma \in \mathcal{F}\) , \(\sigma \in \mathcal{G}^{\mathbb{Z}}(\mathfrak{p}')\) and \(\bar{\sigma} = u\) , which completes the proof of (ii) in this case.

\begin{enumerate}
\def\labelenumi{(\Alph{enumi})}
\setcounter{enumi}{2}
\tightlist
\item
  General case. The field of invariants \(K_{1}\) of G is a radicial extension of K (Algebra, Chapter V, § 10, no. 9, Proposition 14); there therefore exists a single prime ideal \(p_{1}\) of \(A_{1} = A' \cap K_{1}\) lying above p (Lemma 4). If \(p'\) and \(q'\) are two prime ideals of \(A'\) lying above p, then they lie above \(p_{1}\) ; as \(K'\) is a Galois extension of \(K_{1}\) and \(A' \cap K_{1}\) is integrally closed (§ 1, no. 2, Proposition 7 and Corollary to Proposition 8), it follows from (B) that there exists \(\sigma \in G\) such that \(\sigma \cdot p' = q'\) ; whence (i). On the other hand, clearly the field of fractions \(k_{1}\) of \(A_{1}/p_{1}\) is a radicial extension of k (A being integrally closed); as \(k'\) is a quasi-Galois extension of \(k_{1}\) by (B), \(k'\) is a quasi-Galois extension of k, every k-isomorphism of \(k'\) to an algebraically closed extension of \(k'\) being a \(k_{1}\) -isomorphism. This last remark shows also, taking account of (B), that every k-automorphism of \(k'\) is of the form \(\bar{\sigma}\) where \(\sigma \in \mathcal{G}^{\mathbb{Z}}(\mathfrak{p}')\) , which completes the proof of (ii).
\end{enumerate}

Remark

\begin{enumerate}
\def\labelenumi{(\arabic{enumi})}
\setcounter{enumi}{1}
\tightlist
\item
  Suppose that \(\mathbf{K}'\) is a Galois extension of \(\mathbf{K}\) and let us keep the notation of the proof of Proposition 6; for all \(\alpha\) let \(\mathcal{G}_{\alpha}^{\mathbb{Z}}\) (resp. \(\mathcal{G}_{\alpha}^{\mathrm{T}}\) ) be the subgroup of \(\mathcal{G}\) consisting of the \(\sigma\) whose restriction to \(\mathrm{A}' \bigcap \mathrm{K}_{\alpha}\) belongs to \(\mathcal{G}^{\mathbb{Z}}(\mathfrak{p}' \bigcap \mathrm{K}_{\alpha})\) (resp. \(\mathcal{G}^{\mathrm{T}}(\mathfrak{p}' \bigcap \mathrm{K}_{\alpha})\) ). The proof of Proposition 6 shows that these subgroups are closed in \(\mathcal{G}\) and that
\end{enumerate}

\[
\mathcal {G} ^ {\mathrm{Z}} \left(\mathfrak {p} ^ {\prime}\right) = \bigcap_ {\alpha} \mathcal {G} _ {\alpha} ^ {\mathrm{Z}} \quad \text { and } \quad \mathcal {G} ^ {\mathrm{T}} \left(\mathfrak {p} ^ {\prime}\right) = \bigcap_ {\alpha} \mathcal {G} _ {\alpha} ^ {\mathrm{T}}.
\]

Moreover, the set of restrictions to \(A^{\prime} \cap K_{\alpha}\) of the elements of G (resp. \(G_{\alpha}^{T}\) ) is the whole of the group \(\mathcal{G}^{\mathtt{Z}}(\mathfrak{p}^{\prime} \cap \mathtt{K}_{\alpha})\) (resp. \(\mathcal{G}^{\mathtt{T}}(\mathfrak{p}^{\prime} \cap \mathtt{K}_{\alpha})\) ), every K-automorphism of \(K_{\alpha}\) extending to an element of G.

With the same hypotheses, the ring \(\mathbf{A}^{\mathbf{Z}}(\mathfrak{p}')\) (resp. \(\mathbf{A}^{\mathbf{T}}(\mathfrak{p}')\) ) is the union of the directed family of the \(\mathbf{A}^{\mathbf{Z}}(\mathfrak{p}' \cap \mathbf{K}_{\alpha})\) (resp. \(\mathbf{A}^{\mathbf{T}}(\mathfrak{p}' \cap \mathbf{K}_{\alpha})\) ): in fact, every \(x \in \mathbf{A}^{\mathbf{Z}}(\mathfrak{p}')\) (resp. every \(x \in \mathbf{A}^{\mathbf{T}}(\mathfrak{p}')\) ) belongs to one of the \(\mathbf{K}_{\alpha}\) and by the above there exists a \(\beta\) such that \(\mathbf{K}_{\alpha} \subset \mathbf{K}_{\beta}\) and the restrictions to \(\mathbf{A}' \cap \mathbf{K}_{\alpha}\) of the elements of \(\mathcal{G}^{\mathbf{Z}}(\mathfrak{p}')\) (resp. \(\mathcal{G}^{\mathbf{T}}(\mathfrak{p}')\) ) are the same as the restrictions to \(\mathbf{A}' \cap \mathbf{K}_{\alpha}\) of the elements of \(\mathcal{G}^{\mathbf{Z}}(\mathfrak{p}' \cap \mathbf{K}_{\beta})\) (resp. \(\mathcal{G}^{\mathbf{T}}(\mathfrak{p}' \cap \mathbf{K}_{\beta})\) ), the groups \(\mathcal{G}^{\mathbf{Z}}(\mathfrak{p}' \cap \mathbf{K}_{\alpha})\) and \(\mathcal{G}^{\mathbf{T}}(\mathfrak{p}' \cap \mathbf{K}_{\beta})\) being finite; hence \(x\) belongs to \(\mathbf{A}^{\mathbf{Z}}(\mathfrak{p}' \cap \mathbf{K}_{\beta})\) (resp. \(\mathbf{A}^{\mathbf{T}}(\mathfrak{p}' \cap \mathbf{K}_{\beta})\) ).

COROLLARY 1. Under the hypotheses of Proposition 6, let \(f\) be a homomorphism from \(\mathbf{A}\) to a field \(\mathbf{L}\) and \(g_1\) , \(g_2\) two homomorphisms from \(\mathbf{A}'\) to \(\mathbf{L}\) which extend \(f\) . Then there exists a \(\mathbf{K}\) -automorphism \(\sigma\) of \(\mathbf{K}'\) such that \(g_1 = g_2 \circ \sigma\) .

The proof starting from Proposition 6 is the same as that of the Corollary to Theorem 2 starting from the latter.

COROLLARY 2. Let A be an integrally closed domain, K its field of fractions, K' a finite algebraic extension of K and A' a subring of K' containing A and integral over A.

\begin{enumerate}
\def\labelenumi{(\roman{enumi})}
\item
  For every prime ideal \(\mathfrak{p}\) of \(\mathbf{A}\) the set of prime ideals of \(\mathbf{A}'\) lying above \(\mathfrak{p}\) is finite.\\
\item
  If \(\mathfrak{p}'\) is a prime ideal of \(\mathbf{A}'\) lying above \(\mathfrak{p}\) , every element of \(\mathbf{A}' / \mathfrak{p}'\) is of degree \(\leqslant [\mathbf{K}' : \mathbf{K}]\) over the field of fractions of \(\mathbf{A} / \mathfrak{p}\) .
\item
  Let \(\mathbf{K}''\) be the quasi-Galois extension of \(\mathbf{K}\) generated by \(\mathbf{K}'\) in an algebraic closure of \(\mathbf{K}'\) and \(\mathbf{A}''\) the integral closure of \(\mathbf{A}\) in \(\mathbf{K}''\) . The field \(\mathbf{K}''\) is a finite extension of \(\mathbf{K}\) (Algebra, Chapter V, § 6, no. 3, Corollary 1 to Proposition 9) and hence its group of K-automorphisms is finite; it follows that the set of prime ideals of \(\mathbf{A}''\) lying above \(p\) is finite (Proposition 6 (i)). On the other hand, as \(\mathbf{A}''\) is integral over \(\mathbf{A}\) , the mapping \(p'' \mapsto p'' \cap \mathbf{A}'\) of the set of prime ideals of \(\mathbf{A}''\) lying above \(p\) to the set of prime ideals of \(\mathbf{A}'\) lying above \(x\) is surjective (no. 1, Theorem 1).
\item
  The coefficients of the minimal polynomial (over K) of any element \(x' \in A'\) belong to A (§ 1, no. 3, Corollary to Proposition 10); applying the canonical homomorphism \(\pi': A' \to A'/p'\) to the coefficients of this polynomial, an equation of integral dependence with coefficients in A/p and of degree \(\leqslant [K': K]\) is obtained for the class mod. \(p'\) of \(x\) ; whence the conclusion.
\end{enumerate}

COROLLARY 3. With the hypotheses and notation of Corollary 2, if A is semi-local, so is A'.

For every maximal ideal \(m'\) of \(A'\) , \(m' \cap A\) is a maximal ideal of A (no. 1, Proposition 1); the corollary then follows from Corollary 2, since by hypothesis the set of maximal ideals of A is finite.

COROLLARY 4. Let A be an integrally closed domain, K its field of fractions, K' a Galois extension of K, A' the integral closure of A in K', p' a prime ideal of A', p = A ∩ p' and k and k' the fields of fractions of A/p and A'/p' respectively. Then:

\begin{enumerate}
\def\labelenumi{(\roman{enumi})}
\item
  The field of fractions of \(\mathbf{A}^{\mathbf{z}} / (\mathfrak{p}'\cap \mathbf{A}^{\mathbf{z}})\) is equal to \(k\) and the maximal ideal of the local ring of \(\mathbf{A}^{\mathbf{z}}\) relative to \(\mathfrak{p}'\cap \mathbf{A}^{\mathbf{z}}\) is generated by \(\mathfrak{p}\) .
\item
  The field of fractions \(k^{\mathbf{T}}\) of \(\mathbf{A}^{\mathbf{T}} / (\mathfrak{p}' \cap \mathbf{A}^{\mathbf{T}})\) is the greatest separable extension of \(k\) contained in \(k'\) .
\end{enumerate}

The ring A is the ring of invariants in A' of the Galois group of K' over K; if K' is of finite degree over K, the corollary then follows from Propositions 4 and 5 of no. 2. Consider now the general case, K' therefore being the union of a right directed set (K \(_{\alpha}\) ) of finite Galois extensions of K. Then:

\begin{enumerate}
\def\labelenumi{(\roman{enumi})}
\tightlist
\item
  If \(x, y\) are two elements of \(\mathbf{A}^z\) , where \(y \notin \mathfrak{p}'\) , there is an index \(\alpha\) such that \(x\) and \(y\) belong to \(\mathbf{A}^z(\mathfrak{p}' \cap \mathbf{K}_\alpha)\) (Remark 2); by Proposition 4 of no. 2, there are \(x_0, y_0\) in \(\mathbf{A}\) with \(y_0 \notin \mathfrak{p}'\) such that \(xy_0 - x_0y \in \mathfrak{p}'\) , which proves the first assertion of (i); if further \(x \in \mathfrak{p}'\) , we may assume that \(y_0\) satisfies
\end{enumerate}

\[
x y _ {0} \in \mathfrak {p A} ^ {\mathrm{z}} (\mathfrak {p} ^ {\prime} \cap \mathrm{K} _ {\alpha}) \subset \mathfrak {p A} ^ {\mathrm{z}} (\mathfrak {p} ^ {\prime}),
\]

which proves the second assertion of (i).

\begin{enumerate}
\def\labelenumi{(\roman{enumi})}
\setcounter{enumi}{1}
\tightlist
\item
  Suppose now that \(x \in \mathbf{A}^{\mathrm{T}}\) ; there exists \(\alpha\) such that \(x \in \mathbf{A}^{\mathrm{T}}(\mathfrak{p}' \cap \mathbf{K}_{\alpha})\) (Remark 2) and Proposition 5 of no. 2 shows that the class \(\bar{x}\) of \(x \bmod. (\mathfrak{p}' \cap \mathbf{K}_{\alpha} \cap \mathbf{A}^{\mathrm{T}})\) is algebraic and separable over \(k\) ; a fortiori the class mod. \((\mathfrak{p}' \cap \mathbf{A}^{\mathrm{T}})\) of \(x\) is separable over \(k\) ; to complete the proof of the corollary, it is sufficient to show that \(k'\) is a radicial extension of \(k^{\mathrm{T}}\) . Now, \(k'\) is the union of the right directed family of fields of fractions \(k_{\alpha}\) of the rings \((\mathbf{A}' \cap \mathbf{K}_{\alpha}) / (\mathfrak{p}' \cap \mathbf{K}_{\alpha})\) . It follows therefore from Proposition 5 that, if an element of \(k'\) belongs to \(k_{\alpha}\) , it is radicial over the field of fractions of
\end{enumerate}

\[
\mathbf {A} ^ {\mathrm{T}} \left(\mathfrak {p} ^ {\prime} \cap \mathbf {K} _ {\alpha}\right) / \left(\mathfrak {p} ^ {\prime} \cap \mathbf {A} ^ {\mathrm{T}} \left(\mathfrak {p} ^ {\prime} \cap \mathbf {K} _ {\alpha}\right)\right)
\]

and a fortiori over \(k^{\mathbf{T}}\) (by virtue of Remark 2).

DEFINITION 5. With the hypotheses and notation of Proposition 6, the field of invariants \(\mathbf{K}^{\mathrm{Z}}(\mathfrak{p}')\) (resp. \(\mathbf{K}^{\mathrm{T}}(\mathfrak{p}')\) ) of the group \(\mathcal{G}^{\mathrm{Z}}(\mathfrak{p}')\) (resp. \(\mathcal{G}^{\mathrm{T}}(\mathfrak{p}')\) ) in the field \(\mathbf{K}'\) is called the decomposition field (resp. inertia field) of \(\mathfrak{p}'\) with respect to \(\mathbf{K}\) .

We also write \(K^{z}\) (resp. \(K^{T}\) ) in place of \(\mathbf{K}^{\mathbf{Z}}(\mathfrak{p}')\) (resp. \(\mathbf{K}^{\mathrm{T}}(\mathfrak{p}')\) ). It follows from §1, no. 9, Proposition 23 that \(K^{z}\) (resp. \(K^{T}\) ) is the field of fractions of the ring \(A^{z}\) (resp. \(A^{T}\) ); \(A^{z}\) (resp. \(A^{T}\) ) is the integral closure of A in \(K^{z}\) (resp. \(K^{T}\) ).

\textbf{Remarks}\par

\begin{enumerate}
\def\labelenumi{(\arabic{enumi})}
\setcounter{enumi}{2}
\tightlist
\item
  Under the conditions of Corollary 4 of Proposition 6 and assuming that \([\mathbf{K}':\mathbf{K}]\) is finite, the number of distinct prime ideals lying above \(p\) is \([\mathbf{K}^{\mathbf{Z}}:\mathbf{K}]\) , this degree being equal to the index \((\mathcal{G}:\mathcal{G}^{\mathbf{Z}})\) by Galois theory; moreover, it follows from Galois theory that
\end{enumerate}

\[
[ \mathbf {K} ^ {\mathrm{T}}: \mathbf {K} ^ {\mathrm{Z}} ] = (\mathcal {G} ^ {\mathrm{Z}}: \mathcal {G} ^ {\mathrm{T}}) = [ k ^ {\mathrm{T}}: k ].\tag{6}
\]

\begin{enumerate}
\def\labelenumi{(\arabic{enumi})}
\setcounter{enumi}{3}
\tightlist
\item
  Let A be an integrally closed domain, K its field of fractions, \(K'\) a finite algebraic extension of K and \(A'\) the integral closure of A in \(K'\) . Then, for every prime ideal p of A, the number of prime ideals of \(A'\) lying above p is at most \([K':K]_{s}\) (the separable factor of the degree of \(K'\) over K). We may first restrict our attention to the case where \(K'\) is a separable extension of K, for in general \(K'\) is a radicial extension of the greatest separable extension \(K_{0}\) of K contained in \(K'\) , \([K':K]_{s} = [K_{0}:K]\) by definition and, if \(A_{0}\) is the integral closure of A in \(K_{0}\) , the prime ideals of \(A_{0}\) and \(A'\) are in one-to-one correspondence (Lemma 4). Suppose therefore that \(K'\) is separable over K and let N be the Galois extension of K generated by \(K'\) in an algebraic closure of K, G its Galois group, B the integral closure of A in N and P a prime ideal of B lying above p.~Let H be the Galois group of N over \(K'\) and \(G^{z}\) the decomposition group of P; the prime ideals of B lying above p are the s.P where \(s \in G\) (no. 2, Theorem 2) and the relation s.P = s'.P means that \(s' = sg\) where \(g \in G^{z}\) . On the other hand in order that s.P ∩ \(K' = s'\) .P ∩ \(K'\) , it is necessary and sufficient that \(s'\) .P = ts.P, where \(t \in H\) (no. 2, Theorem 2), whence finally \(s' = tsg\) where \(t \in H\) and \(g \in G^{z}\) . The number of prime ideals of A' lying above p is therefore equal to the number of classes of G under the equivalence relation ``there exists \(t \in H\) and \(g \in G^{z}\) such that \(s' = tsg\)'' between s and \(s'\) ; clearly this number is at most equal to the index \((\mathcal{G}:\mathcal{H})\) , the number of right cosets of H in G, and \((\mathcal{G}:\mathcal{H}) = [\mathrm{K}':\mathrm{K}]\) by Galois theory.
\end{enumerate}

PROPOSITION 7. Let A be an integrally closed domain, K its field of fractions, K' a Galois extension of K, G its Galois group, A' the integral closure of A in K', p' a prime ideal of A' and p = A ∩ p'. Finally, let L be a subfield of K' containing K and let p(L) = p' ∩ L.

\begin{enumerate}
\def\labelenumi{(\roman{enumi})}
\item
  The decomposition field (resp. inertia field) of \(\mathfrak{p}'\) with respect to \(\mathbf{L}\) is \(\mathbf{L}(\mathbf{K}^{\mathbf{Z}})\) (resp. \(\mathbf{L}(\mathbf{K}^{\mathbf{T}})\) ); if further \(\mathbf{L}\) is a Galois extension of \(\mathbf{K}\) , the decomposition field of \(\mathfrak{p}(\mathbf{L})\) with respect to \(\mathbf{K}\) is \(\mathbf{L} \cap \mathbf{K}^{\mathbf{Z}}\) .
\item
  If \(\mathbf{L}\) is contained in \(\mathbf{K}^{\mathbf{Z}}\) , \(\mathbf{A} / \mathfrak{p}\) and \((\mathbf{A}' \cap \mathbf{L}) / \mathfrak{p}(\mathbf{L})\) have the same field of fractions and in the local ring \(\mathbf{A}' \cap \mathbf{L}\) corresponding to the prime ideal \(\mathfrak{p}(\mathbf{L})\) , the maximal ideal is generated by \(\mathfrak{p}\) . Conversely, if these two conditions hold and further \(\bigcap_{n \geqslant 0} \mathfrak{p}^{n} \mathbf{A}_{\mathfrak{p}'}' = 0\) , \(\mathbf{L}\) is contained in \(\mathbf{K}^{\mathbf{Z}}\) .
\item
  If \(\mathcal{H}\) is the subgroup of \(\mathcal{G}\) leaving invariant the elements of \(\mathbf{L}\) , clearly the decomposition group (resp. inertia group) of \(\mathfrak{p}'\) with respect to \(\mathbf{L}\) is \(\mathcal{G}^{\mathrm{Z}}\cap \mathcal{H}\) (resp. \(\mathcal{G}^{\mathrm{T}}\cap \mathcal{H}\) ) and the first assertion follows from Galois theory if \(\mathbf{K}'\) is a finite Galois extension of \(\mathbf{K}\) (Algebra, Chapter V, § 10, no. 6, Corollary 1 to Theorem 3); in the general case it follows from the fact that \(\mathbf{A}^{\mathrm{Z}}\) (resp. \(\mathbf{A}^{\mathrm{T}}\) ) is the union of the \(\mathbf{A}^{\mathrm{Z}}(\mathfrak{p}'\cap \mathbf{K}_{\alpha})\) (resp. \(\mathbf{A}^{\mathrm{T}}(\mathfrak{p}'\cap \mathbf{K}_{\alpha})\) ) in the notation of Remark 2: every element \(x\in \mathbf{K}'\) belongs to some \(\mathbf{K}_{\alpha}\) and if it is invariant under \(\mathcal{G}^{\mathrm{Z}}(\mathfrak{p}')\cap \mathcal{H}\) (resp. \(\mathcal{G}^{\mathrm{T}}(\mathfrak{p}')\cap \mathcal{H}\) ) it is also invariant under \(\mathcal{G}^{\mathrm{Z}}(\mathfrak{p}'\cap \mathbf{K}_{\beta})\cap \mathcal{H}\) (resp. \(\mathcal{G}^{\mathrm{T}}(\mathfrak{p}'\cap \mathbf{K}_{\beta})\cap \mathcal{H}\) ) for some suitable \(\beta\) ; hence it belongs by the beginning of the argument to
\end{enumerate}

\[
\mathbf {L} \left(\mathrm{K} ^ {\mathrm{Z}} \left(\mathfrak {p} ^ {\prime} \cap \mathrm{K} _ {\alpha}\right)\right) \subset \mathbf {L} \left(\mathrm{K} ^ {\mathrm{Z}}\right) (\text { resp. } \mathbf {L} \left(\mathrm{K} ^ {\mathrm{T}} \left(\mathfrak {p} ^ {\prime} \cap \mathrm{K} _ {\beta}\right)\right) \subset \mathbf {L} \left(\mathrm{K} ^ {\mathrm{T}}\right)).
\]

Suppose now that L is a Galois extension of K; the restriction to L of every \(\sigma\in G^{z}\) then leaves invariant \(\mathfrak{p}(L)=\mathfrak{p}^{\prime}\cap L\) and hence belongs to the decomposition group of \(\mathfrak{p}(L)\) with respect to K. Conversely, let \(\tau\) be an automorphism of L belonging to this group and let \(\sigma\) be an extension of \(\tau\) to a K-automorphism of \(K^{\prime}\) ; we write \(q^{\prime}=\sigma.p^{\prime}\) . As \(p^{\prime}\) and \(q^{\prime}\) both lie above \(\mathfrak{p}(L)\) , there exists an automorphism \(\rho\in H\) such that \(q^{\prime}=\rho.p^{\prime}\) , whence \(\rho^{-1}\sigma\in G^{z}\) and \(\tau\) is the restriction of \(\rho^{-1}\sigma\) to L; in other words, the decomposition group of \(\mathfrak{p}(L)\) with respect to K is identical with the group of restrictions to L of the automorphisms \(\sigma\in G^{z}\) , which proves the second assertion.

\begin{enumerate}
\def\labelenumi{(\roman{enumi})}
\setcounter{enumi}{1}
\tightlist
\item
  To say that \(L \subset K^{z}\) means that \(H \supset G^{z}\) and the assertions of (ii) are therefore special cases of no. 2, Proposition 4 (ii) and (iii) when \([K':K]\) is finite. In the general case the argument is as in the proof of Proposition 6.
\end{enumerate}

